%================================================================
\chapter{Discussion and Conclusions}
\label{ch:conclusions}
%================================================================

\section{Summary of Contributions}
\label{sec:summary}

Looking back over the four stages, a few things are worth pulling
out rather than just restating.

The spatial baseline settled a question I was not sure about at the
start: whether the graph alone carries enough hydraulic context to
make reconstruction learnable. It does, by a wide margin (a 29-fold
MAE improvement over weighted least squares on Modena). Everything
after that builds on a backbone that already understands the
network's steady state.

The temporal step is the one I would defend hardest. It is not a
small accuracy bump on top of the spatial model; it is the
difference between being able to see replay attacks at all and not.
A replayed reading is, by construction, a real reading from the
same sensor, so nothing in a single snapshot can flag it. Six steps
of history is roughly the least that exposes it, and the nine
stability features get most of the replay signal for free, with no
extra parameters.

The mixture of experts is the least surprising of the four and
still earned its place: one model spreads itself too thin across
five attack signatures, and routing to specialists buys about four
points of aggregate F1 with a router that is right 74\% of the
time.

Self-play is where the actual research question lives. Co-training
a learned attacker against the defender, and then making that
attacker a small population with a diversity loss, produces a
defender that is 5.8 points of F1 and 23\% of MAE better on Modena
than the model it started from. The number I care about more is the
held-out one: the same defender is better on two attack types it
never saw, and the single-attacker version is not. That gap is the
thesis in one sentence.

\section{Limitations}
\label{sec:limitations}

\subsection{Replay Attack Remains Partially Undetected}
\label{sec:limitations_replay}

Replay F1 on Modena is 0.196 for the pretrained model, 0.229 for
self-play single, and 0.165 for the Attacker \moe{}. The degradation
from single to \moe{} attacker is expected: the \moe{} expert 3
(\emph{subtle}) discovered a low-profile strategy that concentrates
on replay windows, making replay attacks in the training distribution
harder and therefore harder for the defender to catch.

More fundamentally, replay attacks are hard because a replayed
pressure reading was a genuine observation. The temporal model can
detect that a reading has stopped evolving (autocorrelation features
help here) but cannot distinguish a genuine steady-state period from
a replay unless it has enough context to know the network's hydraulics
have changed. With a six-step window at one-hour resolution, the
context is limited; a longer window would help.

\subsection{Self-Play Hurts Net1}
\label{sec:limitations_net1}

On Net1 (11 nodes), self-play reduces anomaly F1 from 0.728 to 0.654.
The pretrained model on Net1 already achieves near-perfect F1 on
random, noise, and targeted attacks; the only room for improvement
is stealthy and replay. Self-play adds a noisy adversarial signal that
disrupts the carefully learned decision boundaries for the easy
classes without recovering enough of the hard-class signal to
compensate.

This is consistent with the broader literature on adversarial
training and standard accuracy~\cite{ilyas2019adversarial}: models
trained to be robust to adversarial perturbations often sacrifice
performance on the natural test distribution. For small graphs that
are already well-solved, adversarial fine-tuning is not recommended.

\subsection{Mode Collapse in the Attacker MoE}

The Attacker \moe{} converges to two of four experts. Despite the
diversity and balance losses, experts 0 and 1 receive fewer than
10 windows each at test time. Increasing $\lambda_\text{div}$
or reducing the router's capacity might force more uniform usage,
but the two-mode solution does produce a genuinely useful defender.
Achieving a stable four-mode attacker would require either a larger
perturbation budget (to create more distinct attack strategies) or
a stronger diversity objective.

\subsection{Synthetic Data Only}

All attacks in this thesis are synthetically generated by the
corruption pipeline. No field data with known attack events was
used. This is standard practice in the WDN security
literature (no public labelled dataset of real attacks exists for
the benchmark networks), but it means the models have not been
validated against actual attacker behaviour. In particular, real
attackers may have physical-world knowledge (pipe diameters, valve
schedules) that the synthetic generator does not model.

\subsection{Single-Network Models}

Each model is trained from scratch on a single network. There is no
transfer learning between networks: a model trained on Modena has
essentially zero F1 on Net3 without retraining. In a practical
deployment covering a large water utility with many distinct
networks, per-network training is feasible but expensive.
Meta-learning or graph-level transfer approaches are natural
extensions.

\section{Future Work}
\label{sec:future}

\textbf{Longer temporal context.} The six-step window was chosen to
keep the temporal model tractable and the training data size manageable.
A window of 24 steps (full daily cycle) would provide enough context
to detect replay attacks by comparing against the diurnal pattern,
at the cost of increased model size and training time. Transformer-based
temporal models that can attend over longer sequences are worth
exploring.

\textbf{Physics-informed attacker constraints.} The current attacker
is penalised for violating flow conservation but is not constrained
to satisfy it exactly. A gradient-projected attacker that solves a
linearised hydraulic feasibility problem at each step would produce
physically realistic attacks, making the defender stronger against
real adversaries who understand the network's physics.

\textbf{Cross-network transfer.} A model that generalises across
networks would reduce the per-network training cost substantially.
Graph-level features (spectral properties, degree distribution,
pressure distribution summary statistics) could be used to condition
the model on the target network.

\textbf{Online adaptation.} The current models are trained offline
and evaluated on a fixed test set. A practical deployment would
need to adapt to network configuration changes (new pipes, changed
demand patterns) and to novel attack strategies as they emerge. An
online variant of the self-play loop that continuously updates the
defender against new attacker discoveries is a natural direction.

\textbf{Explainability for operators.} The \textsc{GNNExplainer}
integration in the codebase (\texttt{src/wdn/explainability.py})
identifies which neighbouring nodes and edges contributed most to
each anomaly flag. Making this explanation accessible to operators
through the dashboard (as in the Explainability page) is already
implemented, but a formal user study to determine whether the
explanations are actionable is outstanding.

\section{Conclusion}
\label{sec:conclusion}

This thesis showed that graph neural networks, combined with temporal
modelling, mixture of experts, and adversarial self-play, produce a
water distribution network anomaly detector that is both accurate
and robustly generalising. The key insight is that a defender trained
against a learned, co-evolving attacker population is better prepared
for novel attack strategies than a defender trained against a fixed
library of hand-crafted attacks.

The unexpected result was the emergent attack vocabulary. An Attacker
\moe{} with no attack-class supervision split its perturbation space
into two distinct strategies: one that hits hard and broadly (the
\emph{bold} expert) and one that stays quiet, drifting toward
replay-like windows (the \emph{subtle} expert). That split came from
the diversity loss alone, which pushed the experts apart in perturbation
space. A single attacker would not have found both.

Replay attacks remain the Achilles heel of the whole framework.
A reading that was true one hour ago is hard to distinguish from a
reading that is true now, especially under a six-step window. Closing
this gap is the single most important direction for future work.

The complete codebase, trained model weights, and an interactive
Streamlit dashboard are available at the project repository. All
experiments can be reproduced in under four hours on a laptop with
Apple Silicon.
